\documentclass[11pt,a4paper]{article}
\usepackage[utf8]{inputenc}
\usepackage[T1]{fontenc}
\usepackage[brazilian]{babel}
\usepackage{lmodern}
\usepackage{amsmath,amssymb,amsthm}
\usepackage[margin=2.5cm]{geometry}
\usepackage{algorithm}
\usepackage[noend]{algpseudocode}

\floatname{algorithm}{Algoritmo}
\renewcommand{\algorithmicrequire}{\textbf{Entrada:}}
\renewcommand{\algorithmicensure}{\textbf{Saída:}}

\newtheorem{lema}{Lema}
\newtheorem{teorema}{Teorema}

\title{Avaliação Continuada 04 -- Questão 2\\ \large Paradas em hotéis}
\date{}

\begin{document}
\maketitle

\noindent\textbf{Notação.} Uma sequência de hotéis $S=[s_1,\dots,s_r]$ é
\emph{viável} se $s_1 = 1$, $s_r = n$, $s_i < s_{i+1}$ e
$x[s_{i+1}] - x[s_i] \le D$ para todo $1 \le i < r$. A escolha gulosa a partir
da parada $u$ é $\max\{\, w > u : x[w] - x[u] \le D \,\}$, que existe para
$u<n$, pois $x[u+1]-x[u]\le D$.

\noindent\textbf{Exemplo} ($D=12$): a partir do hotel 1 ($x=0$), o mais distante
alcançável é o hotel 4 ($x=11$); a partir do hotel 4, é o hotel 6 ($x=17$).
Logo $G=[1,4,6]$.

%--------------------------------------------------------------------
\section*{(a) Propriedade de escolha gulosa (argumento de troca)}

Sejam $G=[u_1,\dots,u_m]$ e $O=[v_1,\dots,v_\ell]$ como no enunciado, com
$v_i = u_i$ para $1\le i<k$, $v_k \ne u_k$ e $1<k<\ell$.

\begin{lema}
$v_k < u_k$.
\end{lema}
\begin{proof}
Como $O$ é viável e $v_{k-1}=u_{k-1}$, temos $x[v_k] - x[u_{k-1}] \le D$ e
$v_k > u_{k-1}$. Logo $v_k$ é um candidato da escolha gulosa a partir de $u_{k-1}$,
e $u_k$ é o maior candidato: $v_k \le u_k$. Como $v_k \ne u_k$, $v_k < u_k$.
\end{proof}

\noindent\textbf{Construção.} Troque $v_k$ por $u_k$ em $O$, descartando os
hotéis de $O$ que ficaram ``para trás'' de $u_k$:
\[
O' = [\,v_1,\dots,v_{k-1},\; u_k,\; v_{j+1},\dots,v_\ell\,],
\qquad j = \max\{\, i \ge k : v_i \le u_k \,\}.
\]
(Note que $j\ge k$, pois $v_k<u_k$. Se $v_j = u_k$, então $v_j$ e $u_k$ são o mesmo
hotel e ele aparece uma única vez.)

\begin{lema}
$O'$ é viável e $|O'| \le |O|$.
\end{lema}
\begin{proof}
\emph{Tamanho:} $O'$ remove de $O$ os $j-k+1 \ge 1$ hotéis $v_k,\dots,v_j$ e
insere apenas $u_k$, logo $|O'| = \ell - (j-k) \le \ell$.

\emph{Ordem:} $v_{k-1}=u_{k-1}<u_k$ e, pela escolha de $j$, $u_k < v_{j+1}$.

\emph{Distâncias:} os trechos $v_1\to\cdots\to v_{k-1}$ e
$v_{j+1}\to\cdots\to v_\ell$ são trechos de $O$, portanto $\le D$. Restam dois trechos:
\begin{itemize}
  \item $u_{k-1}\to u_k$: $x[u_k]-x[u_{k-1}] \le D$ pela definição da escolha gulosa;
  \item $u_k \to v_{j+1}$ (se $j<\ell$): como $x[v_j] \le x[u_k]$,
        \[ x[v_{j+1}] - x[u_k] \;\le\; x[v_{j+1}] - x[v_j] \;\le\; D. \]
\end{itemize}
\emph{Extremos:} $O'$ começa em $v_1=1$. Se $j<\ell$, termina em $v_\ell=n$;
se $j=\ell$, então $n = v_\ell \le u_k \le n$, logo $u_k=n$ e $O'$ termina em $n$.
\end{proof}

Como $O$ é ótima e $|O'|\le|O|$, $O'$ também é ótima (logo, na verdade, $j=k$ e
$|O'|=|O|$), e $O'$ coincide com $G$ nas $k$ primeiras paradas. Isso prova a
sentença: é possível trocar $v_k$ por $u_k$ sem piorar a solução ótima.

\begin{teorema}
A solução gulosa $G$ é ótima, ou seja, $m = \ell$.
\end{teorema}
\begin{proof}
Aplicando a troca repetidamente, o primeiro índice de divergência cresce a cada
passo; obtemos assim uma solução ótima $O^*=[v_1,\dots,v_\ell]$ com $v_i=u_i$
para todo $1\le i\le \ell-1$. Como $O^*$ é viável, $x[n]-x[u_{\ell-1}] =
x[v_\ell]-x[v_{\ell-1}] \le D$; então o hotel $n$ (o mais distante de todos) é
alcançável a partir de $u_{\ell-1}$ e a escolha gulosa é $u_\ell = n$. Assim o
guloso termina com $m=\ell$ paradas. (Equivalentemente, pelo subproblema descrito no
enunciado, a cada troca o problema restante parte de $x[u_k]$ e o argumento se
repete por indução.)
\end{proof}

%--------------------------------------------------------------------
\section*{(b) Pseudocódigo e complexidade}

\begin{algorithm}[H]
\caption{\textsc{Paradas-Hotéis}$(x, n, D)$}
\begin{algorithmic}[1]
\Require posições $x[1..n]$ em ordem crescente, com $x[i+1]-x[i]\le D$
\Ensure lista $G$ com o menor número de paradas, de $1$ até $n$
\State $G \gets [1]$
\State $u \gets 1$ \Comment{última parada}
\State $i \gets 1$ \Comment{candidato atual}
\While{$u < n$}
  \While{$i < n$ \textbf{and} $x[i+1] - x[u] \le D$}
    \State $i \gets i + 1$ \Comment{avança enquanto o próximo hotel é alcançável}
  \EndWhile
  \State $u \gets i$ \Comment{escolha gulosa: hotel mais distante alcançável}
  \State \textsc{Anexa}$(G, u)$
\EndWhile
\State \Return $G$
\end{algorithmic}
\end{algorithm}

\noindent\textbf{Terminação e corretude do laço.} No início de cada iteração do
laço externo, $i = u$. O laço interno avança $i$ ao maior índice com
$x[i]-x[u]\le D$ (a escolha gulosa); como $x[u+1]-x[u]\le D$, ele avança pelo menos
uma vez, logo $u$ cresce estritamente e o algoritmo termina com $u=n$.

\noindent\textbf{Complexidade.} O índice $i$ nunca decresce e vai de $1$ até $n$,
então o laço interno executa no total no máximo $n-1$ incrementos ao longo de toda
a execução. O laço externo executa no máximo $n-1$ vezes, cada uma com custo
$O(1)$ fora do laço interno (anexar a $G$ custa $O(1)$). Logo o tempo total é
\[ T(n) = O(n), \]
sem necessidade de ordenação, pois $x$ já está em ordem crescente. O espaço adicional
é $O(n)$ para armazenar $G$ (ou $O(1)$, se as paradas forem apenas impressas).

\end{document}
